\documentclass{article}
\usepackage{amsmath}
\usepackage{graphicx}
\usepackage{hyperref}

\title{A Clean LaTeX Document}
\author{Author Name}
\date{\today}

\begin{document}

\maketitle

\section{Introduction}

This is a clean LaTeX document with no issues. It compiles correctly and follows good typesetting practices.

As shown in Figure~\ref{fig:example}, the results are consistent with our hypothesis. Previous work~\cite{smith2020} also supports this finding.

\section{Methods}

We used the standard approach described in~\cite{jones2019}.

\begin{figure}[htbp]
  \centering
  \includegraphics[width=0.8\textwidth]{example.png}
  \caption{An example figure.}
  \label{fig:example}
\end{figure}

\section{Results}

The results are shown in Table~\ref{tab:results}.

\begin{table}[htbp]
  \centering
  \begin{tabular}{lcc}
    \hline
    Method & Accuracy & Time \\
    \hline
    A & 0.95 & 10s \\
    B & 0.92 & 5s \\
    \hline
  \end{tabular}
  \caption{Results comparison.}
  \label{tab:results}
\end{table}

\section{Conclusion}

In conclusion, our method achieves the best results.

\bibliographystyle{plain}
\bibliography{references}

\end{document}
